\documentclass[11pt]{article}
\usepackage[a4paper,margin=25mm]{geometry}
\usepackage{amsmath,amssymb,amsthm,mathtools}
\usepackage[T1]{fontenc}
\usepackage{lmodern,microtype}
\usepackage{hyperref}
\hypersetup{colorlinks=true,linkcolor=blue,citecolor=blue,urlcolor=blue}
\newcommand{\norm}[1]{\left\lVert#1\right\rVert}
\newcommand{\ip}[2]{\langle#1,#2\rangle}
\newcommand{\zer}{\operatorname{zer}}
\newcommand{\ran}{\operatorname{ran}}
\newcommand{\gph}{\operatorname{gph}}
\newcommand{\dist}{\operatorname{dist}}
\newtheorem{theorem}{Theorem}[section]
\newtheorem{proposition}[theorem]{Proposition}
\newtheorem{lemma}[theorem]{Lemma}
\newtheorem{corollary}[theorem]{Corollary}
\theoremstyle{remark}
\newtheorem{remark}[theorem]{Remark}
\title{General-modulus convergence of generalized\ proximal point iterations}
\author{PPA research project}
\date{October 9, 2026}
\begin{document}
\maketitle
\begin{abstract}
We transfer reflected graph moduli and true residual error bounds to the
generalized proximal point update
$v(x_k)-v(x_{k+1})\in\lambda F(x_{k+1})$.
A full-fiber kernel-coordinate relation yields direct, residual-energy,
and anchored-dissipation scalar certificates. The energy branch contains
every adaptive-strong-monotonicity instance of the fixed 2026 GPPA
Theorem 2, including arbitrarily small positive parameters. General
moduli add nonlinear and nonpower summability mechanisms. Physical point
convergence and physical finite length require distinct lifting conditions.
We give complete two-branch examples with nonlinear strongly monotone
kernels and any prescribed physical order greater than one, a full-graph
pair-monotone representation obstruction, and a sharp inverse-H\"older
spiral boundary. Persistent output errors yield a nonlinear tube, rather
than unconditional exact convergence. This is an internally reconstructed
research note; global originality and unrestricted reformulation
comparisons remain separate obligations.
\end{abstract}

\section{The complete kernel-coordinate relation}
Let $H$ be a real Hilbert space, $F:H\rightrightarrows H$ a complete
relation, $v:H\to H$ an everywhere-defined single-valued kernel, and
$\lambda>0$. The kernel itself need not be monotone or injective. Define
\begin{equation}\label{eq:union}
 Q=\ran v,\quad S=\zer F,\quad
 A(z)=\bigcup_{v(x)=z}F(x),\quad
 r_A(z)=\inf_{f\in A(z)}\norm f.
\end{equation}
Empty fibers have residual $+\infty$. Directly,
\begin{equation}\label{eq:zeros}
 \zer A=v(S),\qquad r_A(v(x))\le r_F(x).
\end{equation}
The inequality can be strict for a noninjective kernel. For a physical
graph block $\mathcal G\subset\gph F$, put
$\Gamma=\{(v(x),f):(x,f)\in\mathcal G\}$.
Then a block GPPA step is exactly
\begin{equation}\label{eq:step}
 z_k=z_{k+1}+\lambda f_{k+1},\quad
 (z_{k+1},f_{k+1})\in\Gamma,\quad z_k=v(x_k).
\end{equation}
Every kernel output has a physical representative by the definition of
$\Gamma$. For the full graph $\mathcal G=\gph F$, no branch or residual
fiber is deleted. Existence only needs to be required on $Q$, not on all
of $H$. Neither completeness nor closedness of $Q$ follows automatically.

\section{A combined scalar convergence theorem}
Fix a nonempty closed $Z\subset\zer A$, an open $U\subset H$, and $R>0$.
Set $E=\{z\in Q\cap U:d(z,Z)<R\}$. Assume the following same-object data.
\begin{enumerate}
\item Every $z\in E$ has an output on $\Gamma$, and $(p,0)\in\Gamma$
for every $p\in Z$.
\item A continuous nondecreasing $\omega:[0,R]\to[0,\infty)$,
$\omega(0)=0$, satisfies, on every pair in $\Gamma$ with input distance
$\delta\le R$,
\begin{equation}\label{eq:rl}
 \norm{\Delta y-\lambda\Delta f}
 \le\omega(\norm{\Delta y+\lambda\Delta f}).
\end{equation}
\item For all actual outputs with $r_A(y)\le\bar t$, a continuous
strictly increasing gauge $\psi:[0,\bar t]\to[0,\infty)$, $\psi(0)=0$,
gives $d(y,Z)\le\psi(r_A(y))$. Put
\begin{equation}\label{eq:bc}
 B(t)=\frac{t+\omega(t)}2,\quad
 C(t)=\frac{t^2+\omega(t)^2}2,\qquad B(R)/\lambda\le\bar t.
\end{equation}
\end{enumerate}
An optional additional condition is a continuous nondecreasing
$a:[0,M]\to[0,\infty)$, $a(0)=0$, such that
\begin{equation}\label{eq:anchor}
 \ip f{y-p}\ge a(d(y,Z))\quad(p\in Z)
\end{equation}
at each actual output. Here $M=\psi(\bar t)$ and $g=\psi^{-1}$ on $[0,M]$.

\begin{lemma}\label{lem:one}
For $d=d(z,Z)$, $d^+=d(z^+,Z)$ and $s=\norm{z-z^+}$, each covered step
satisfies
\begin{align}\label{eq:one}
 (d^+)^2+s^2&\le C(d),&s&\le B(d),\\
 r_A(z^+)&\le s/\lambda\le\bar t,&
 d^+&\le\psi(s/\lambda),&s&\ge\lambda g(d^+).\nonumber
\end{align}
Under \eqref{eq:anchor} it also satisfies
\begin{equation}\label{eq:oneanchor}
 (d^+)^2+s^2+2\lambda a(d^+)\le d^2.
\end{equation}
The kernel output and its selected graph value are unique; the physical
output can remain multivalued.
\end{lemma}
\begin{proof}
Choose $p_n\in Z$ with $\norm{z-p_n}\to d<R$. For large $n$ these
anchors are inside the comparison cutoff. Write $b=z-z^+=\lambda f$
and $h_n=z^+-p_n$. Then $h_n+b=z-p_n$ and
$\norm{h_n-b}\le\omega(\norm{z-p_n})$. The parallelogram identity and
the bounds on $2b=(h_n+b)-(h_n-b)$ give the first line of
\eqref{eq:one} after taking the limit. The selected graph value belongs
to the full fiber $A(z^+)$; this opens the residual window before the EB
is used. Monotonicity of $\psi$, followed by its inverse, gives the
second line. Polarization and \eqref{eq:anchor} give
\eqref{eq:oneanchor}. Same-input comparison has $\delta=0$ and forces
both $\Delta y+\lambda\Delta f$ and $\Delta y-\lambda\Delta f$ to
vanish. At $d=0$, comparison with $(z,0)$ forces $s=0$ and $f=0$.
No Hilbert-space nearest-point attainment was assumed.
\end{proof}

Define, on $[0,M]$,
\begin{equation}
 V(t)=t^2+\lambda^2g(t)^2,\qquad
 W(t)=V(t)+2\lambda a(t).
\end{equation}
For either increasing function $J=V,W$, write $J^{[-1]}(u)=J^{-1}(u)$
when $u\le J(M)$, and $J^{[-1]}(u)=M$ otherwise. The cap is justified
by the already proved $d^+\le M$. Define
\begin{equation}\label{eq:tau}
 \tau(t)=\min\left\{\psi(B(t)/\lambda),\ V^{[-1]}(C(t))\right\},
\end{equation}
adding $W^{[-1]}(t^2)$ to the minimum when \eqref{eq:anchor} is imposed.
Lemma~\ref{lem:one} gives $d^+\le\tau(d)$, because
\begin{equation}\label{eq:energy}
 V(d^+)\le C(d),\qquad W(d^+)\le d^2\quad\text{when available}.
\end{equation}

\begin{theorem}\label{thm:main}
Assume the displayed data and $\tau(t)<t$ for $0<t\le R$. For an
initial $z_0=v(x_0)\in E$, put $d_0=d(z_0,Z)$ and
\begin{equation}\label{eq:budget}
 H_\tau(r)=\sum_{j\ge0}B(\tau^j(r)).
\end{equation}
If $H_\tau(d_0)<\infty$ and
$H_\tau(d_0)<d(z_0,H\setminus U)$, all block GPPA choices continue
indefinitely. Their unique kernel trajectory stays in $E$, has finite
length, and converges to $z_\infty\in Z\cap U$. Moreover
\begin{equation}
 d(z_k,Z)\le\tau^k(d_0),\qquad
 \norm{z_k-z_\infty}\le H_\tau(\tau^k(d_0)).
\end{equation}
\end{theorem}
\begin{proof}
$\tau$ is continuous nondecreasing and vanishes at zero. Its iterates
decrease to zero, since a positive limit would be a fixed point.
Every valid finite prefix has step bounds $B(\tau^j(d_0))$ and remains
within the strict initial budget of $z_0$. Its distance to $Z$ remains
below $R$, and its next output is in $Q$ by construction. This proves
retention and permits the next use of coverage without circularity.
Summable steps give a Cauchy sequence in the ambient Hilbert space.
Distance to closed $Z$ tends to zero, so the limit belongs to $Z$ and
hence to $Q$. The strict budget retains the limit in $U$. The tail of
\eqref{eq:budget} proves the stated estimate.
\end{proof}

Two useful independent sufficient certificates are as follows.
If $\tau(t)\le\kappa t$, $0<\kappa<1$, and
\begin{equation}\label{eq:dini}
 \int_0^R\omega(t)\,\frac{dt}{t}<\infty,
\end{equation}
then $H_\tau(r)\le\tfrac12\sum_{j\ge0}[\kappa^jr+\omega(\kappa^jr)]$.
Integral comparison on $[\kappa^{j+1}r,\kappa^jr]$, for $r>0$, proves
equivalence of Dini integrability and geometric sampling of $\omega$.
It does not assert necessity for the possibly faster $H_\tau$ series.
Alternatively, if $R\le M$ and $C(t)\le qV(t)$, $q<1$, the energy
branch gives $V(d_{k+1})\le qV(d_k)$ and
$s_k\le q^{(k+1)/2}\sqrt{V(d_0)}$. Its separate total-length budget is
\begin{equation}\label{eq:energybudget}
 \frac{\sqrt{qV(d_0)}}{1-\sqrt q}.
\end{equation}
The same strict retention proof applies with this budget, without
requiring \eqref{eq:dini} or summability of a looser $B$ envelope.

\section{Inclusion of the fixed GPPA quadratic theory}
The source \cite{LMT} requires adaptive strong monotonicity (ASM)
\begin{equation}\label{eq:asm}
 \ip{f-f'}{v(x)-v(x')}\ge\varepsilon\norm{v(x)-v(x')}^2
 \quad((x,f),(x',f')\in\gph F),
\end{equation}
nonempty $S$, and $Q\subset\ran(v+\lambda F)$. These assumptions do
not require $v$ itself to be monotone, injective, or surjective.
Two zero graph points in \eqref{eq:asm} give $v(S)=\{z_*\}$. The full
relation $A$ is monotone, so $\omega(t)=t$ is valid. Zero comparison
and Cauchy--Schwarz give, for every full graph value,
$\norm{z-z_*}\le\norm f/\varepsilon$. Taking the infimum yields
the true kernel EB $\psi(t)=t/\varepsilon$.
Consequently
\begin{equation}\label{eq:asminclusion}
 V(t)=(1+\lambda^2\varepsilon^2)t^2,\qquad
 \tau_E(t)=\frac{t}{\sqrt{1+\lambda^2\varepsilon^2}}<t.
\end{equation}
Thus the energy alternative contains every source Theorem~2 kernel
case for every $\lambda\varepsilon>0$. The direct test alone would
require $\lambda\varepsilon>1$ with $\omega=t$, and is insufficient
to establish inclusion. Even with the sharp modulus for $F=\varepsilon I$
and $v=I$, at $\lambda\varepsilon=1/4$ the direct ratio is at least
$16/5$, while the energy branch contracts.

Keeping the source anchor dissipation $a(t)=\varepsilon t^2$ gives
\begin{equation}\label{eq:asmrate}
 W(t)=(1+\lambda\varepsilon)^2t^2,\qquad
 \norm{z_k-z_*}\le(1+\lambda\varepsilon)^{-k}\norm{z_0-z_*}.
\end{equation}
This elementary consequence of the source hypotheses sharpens its
printed factor $(1+2\lambda\varepsilon)^{-1/2}$; it is not claimed
as an independently novel mechanism. Its physical R-continuity
assumption transfers selected residuals to
$d(x_{k+1},S)\le\rho(s_k/\lambda)$ once the selected residual enters
its stated local radius; geometric kernel decay ensures eventual entry.
A global inverse bound gives this at every step. A linear $\rho$ preserves a
geometric set-distance rate. Neither this bound nor ASM controls
unrestricted movement within a noninjective kernel fiber.

There is also a weak-EB extension with no uniform geometric factor.
For pair-monotone $A$, suppose $0<d_0\le M$ and put $r_j=2^{-j}d_0$.
The case $d_0=0$ is absorbed at zero and needs no shell division.
The descent $d_k^2-d_{k+1}^2\ge s_k^2$ and
$s_k\ge\lambda g(d_{k+1})$ give the shell-length bound
\begin{equation}\label{eq:shell}
 \sum_{j\ge0}\left(r_j+\frac{r_j^2}{\lambda g(r_j/2)}\right).
\end{equation}
Indeed, interior steps of one shell have length at most their squared
length divided by $\lambda g(r_j/2)$; the squared lengths telescope
to at most $r_j^2$. There is at most one crossing step, of length at
most $r_j$. Skipped shells are not charged. A finite strict budget
in \eqref{eq:shell} gives the same continuation and limit proof.
For $\psi(t)=Kt^q$ it suffices that $q>1/2$.
The full scalar example $F(x)=\operatorname{sgn}(x)|x|^p$, $v=I$,
$1<p<2$, has $q=1/p$, full coverage, and polynomial PPA convergence
$x_k\sim[(p-1)\lambda k]^{-1/(p-1)}$ from $x_0>0$; it has no ASM
constant for the same identity kernel. Source Theorem~1 and alternative
kernels remain possible, so this witness concerns the fixed-kernel
Theorem~2 assumption class.

\section{Physical lifting and retractions}
Suppose independently that the actual physical points are controlled by
a continuous nondecreasing inverse or selected-fiber modulus $\eta$:
\begin{equation}\label{eq:lift}
 \norm{x-y}\le\eta(\norm{v(x)-v(y)}),\qquad\eta(0)=0.
\end{equation}
For all-pairs lifting, this bound forces injectivity on its stated
physical set. Noninjective kernels can use a specified section with an
all-endpoint bound; conclusions then belong to that section policy.
If the bound is imposed only on consecutive allowed transitions, it gives
physical Cauchy convergence only after summability of the transformed
step bounds is proved. Continuity of a transition-only modulus is insufficient.
A uniform inverse estimate is required on an output collar before local
retention is proved, not merely on two already-retained inputs.

Kernel finite length and the all-endpoint version of \eqref{eq:lift} imply physical Cauchy convergence
by endpoint tail comparison. Zero membership needs either a continuous
inverse on $Z=v(S)$, or a closed physical zero target and its true EB,
or the appropriate graph-closure implication at residual zero. Physical
finite length is stronger: for the geometric-distance case a sufficient
transformed-envelope certificate, for $r>0$, is
\begin{equation}\label{eq:physicaldini}
 \sum_{j\ge0}\eta(B(\kappa^jr))<\infty
 \quad\Longleftrightarrow\quad
 \int_0^R\eta(B(t))\,\frac{dt}{t}<\infty.
\end{equation}
If $v$ is $m$-strongly monotone, Cauchy--Schwarz supplies
$\eta(t)=t/m$; continuity, coverage, and surjectivity remain separate.
If $\eta(t)=Mt^\alpha$ and $\omega(t)\asymp[\log(e/t)]^{-a}$,
kernel Dini needs $a>1$, whereas \eqref{eq:physicaldini} needs
$a\alpha>1$.

If a homeomorphic kernel conjugates a continuous kernel iteration on
an invariant open domain $\mathcal O_z$, its kernel scalar length series
converges uniformly from a fixed radius, and the strict local budget
retains every limit in the zero target $Z\cap\mathcal O_z$, with all
points of that target fixed by the iteration, then
the finite iterates converge uniformly to a continuous kernel retraction
$\Pi_z$. Its physical counterpart is
$\Pi_x=v^{-1}\Pi_zv$. This endpoint result only needs kernel Dini;
physical finite length need not follow.

\section{Nonlinear kernels and arbitrary physical orders}
Let $G$ be a full relation and $V$ a homeomorphism. Set
\begin{equation}\label{eq:conjugate}
 F^V(x)=G(Vx),\qquad
 J_{\lambda F^V}^V=V^{-1}J_{\lambda G}V.
\end{equation}
Graph values are left unchanged. All complete fibers, kernel RL moduli,
and true residuals transfer exactly, with
$r_{F^V}(x)=r_G(Vx)$.
A concrete genuinely nonlinear strongly monotone kernel in $\mathbb R^2$ is
\begin{equation}\label{eq:nonlinearkernel}
 V(\xi,y)=(\xi+c\sin y,\ y+b\tanh y),\quad b>0,\quad0<c<1.
\end{equation}
It is globally invertible. Writing $\varphi=(y+b\tanh y)^{-1}$,
its inverse is $(p,r)\mapsto(p-c\sin\varphi(r),\varphi(r))$.
Its strong monotonicity constant is at least $1-c/2$; its inverse
Lipschitz constant is at most $1+c$.

Fix $0<\gamma<1$, $\nu>1$ and $A_0>0$. Define the full closed graph
\begin{equation}\label{eq:doublegraph}
 G(\xi,Y)=
 \{(-A_0Y^{\gamma/\nu},Y^{1/\nu}-Y),
    (-A_0Y^{\gamma/\nu},-Y^{1/\nu}-Y)\}\quad(Y\ge0),
\end{equation}
and $G(\xi,Y)=\varnothing$ for $Y<0$. Its full step-one resolvent is
\begin{equation}\label{eq:doubleprox}
 J_G(p,r)=(p+A_0|r|^\gamma,|r|^\nu),\qquad S_G=\mathbb R\times\{0\}.
\end{equation}
On input heights $|r|\le R_0<1$ and input comparison distances $\delta\le R$,
the corresponding full two-branch graph block has reflected modulus
\begin{equation}
 \omega_R(t)=L_Rt^\gamma,\quad
 L_R=2A_0+(1+2\nu R_0^{\nu-1})R^{1-\gamma}.
\end{equation}
This follows from $||r|^\gamma-|s|^\gamma|\le|r-s|^\gamma$
and the derivative bound for $2|r|^\nu-r$ on the collar. All cross-branch
pairs are included. Its full residual is
\begin{equation}
 r_G(\xi,Y)^2=A_0^2Y^{2\gamma/\nu}+(Y^{1/\nu}-Y)^2,
 \qquad Y\le(r_G(\xi,Y)/A_0)^{\nu/\gamma}.
\end{equation}
The smaller positive branch determines the true infimum even at a
negative input selecting the other branch. For sufficiently small $R$,
the direct compatibility ratio is bounded by
\begin{equation}
 R^{\nu-1}\left(\frac{R^{1-\gamma}+L_R}{2A_0}\right)^{\nu/\gamma}<1.
\end{equation}
Full coverage on the collar, all zero anchors, Dini, and the residual
evaluation window hold for this same graph block; the normal update
independently proves collar invariance. Conjugation by
\eqref{eq:nonlinearkernel} gives the corresponding complete physical
generalized iteration.

From $0<r_0<1$, the exact normal coordinate is $r_k=r_0^{\nu^k}$.
Let $E(r)=\sum_{j\ge0}r^{\gamma\nu^j}$. Since
$\nu^j\ge1+j(\nu-1)$,
\begin{equation}
 r^\gamma\le E(r)\le
 \frac{r^\gamma}{1-r^{\gamma(\nu-1)}},\qquad E(r)\sim r^\gamma.
\end{equation}
The physical point error to its selected zero is
\begin{equation}
 e_k^2=[A_0E(r_k)+c\sin\varphi(r_k)]^2+\varphi(r_k)^2.
\end{equation}
Because $\varphi(r)\sim r/(1+b)$, this gives
\begin{equation}\label{eq:qorder}
 e_k\sim A_0r_k^\gamma,\qquad
 \lim_k\frac{e_{k+1}}{e_k^\nu}=A_0^{1-\nu}.
\end{equation}
Thus every prescribed physical order $\nu>1$, including $\nu>2$, is
realized. This is an existential family, not a universal rate guarantee.

\begin{proposition}[Full-graph representation obstruction]\label{prop:obstruction}
No pair-monotone source kernel can preserve the specified nonstationary
conjugated trajectory of \eqref{eq:doublegraph} on that same full graph,
even with arbitrary positive step sizes and branch choices.
\end{proposition}
\begin{proof}
Any such kernel $K$ would give $Q=K\circ V^{-1}$ pair-monotone on
$G$. On a tangential product slice and a positive connected normal
interval $I\subset(0,1)$, write the two values as
$(-C_0\phi(y),a_+(y))$ and $(-C_0\phi(y),a_-(y))$,
where $\phi$ is strictly increasing locally Lipschitz and
$a_+>0>a_-$. Same-height cross comparisons force the normal component
of $Q$ to be independent of the tangential coordinate. At a fixed
tangential position write its components as $q_1(y),q_2(y)$.
Cross-branch comparisons at $y,z$ give
\begin{equation}
 A\Delta q_2\ge C_0\Delta\phi\,\Delta q_1,\qquad
 -B\Delta q_2\ge C_0\Delta\phi\,\Delta q_1,
\end{equation}
where $A=a_+(y)-a_-(z)>0$, $B=a_+(z)-a_-(y)>0$.
Multiplication by $B,A$ and addition show that $q_1$ is nonincreasing.
On a compact $[\alpha,\beta]\subset I$, the coefficients have a
positive lower bound $m$ and $\phi$ a Lipschitz constant $L$. Hence
\begin{equation}
 |\Delta q_2|\le\frac{C_0L}{m}|y-z|\,|\Delta q_1|.
\end{equation}
Apply this on an equal $N$-part partition. The absolute increments
of $q_1$ telescope to its finite total variation, even if it jumps,
so $|q_2(\beta)-q_2(\alpha)|\le
(C_0L/m)(\beta-\alpha)[q_1(\alpha)-q_1(\beta)]/N\to0$.
Thus the normal kernel component is constant. Along the proposed
trajectory its consecutive difference is zero, whereas both allowed
normal graph values at a positive height are nonzero. This contradicts
any positive-step source update.
\end{proof}
This obstruction concerns the same complete relation and specified
algorithm. Zero-set-preserving positive-branch restrictions can still
import the same trajectory's convergence, as already established in
the project. The obstruction is not a proof of convergence
nonderivability through every legitimate reformulation.

For general power data $\omega(t)=Lt^\gamma$, $\psi(s)=Cs^q$,
$0<\gamma<1$, the direct ratio is exactly
\begin{equation}
 C\left(\frac{L+t^{1-\gamma}}{2\lambda}\right)^q t^{q\gamma-1}.
\end{equation}
It contracts on a small window when $q\gamma>1$, or at $q\gamma=1$
when $C(L/(2\lambda))^q<1$. At critical coefficient one it exceeds
one for every positive $t$. For $\gamma=1$, the coefficient uses
$1+L$, not $L$. A failed sufficient test does not imply divergence.
The source's quadratic energy expression does not mean that its
algorithm is limited to quadratic convergence order.

\section{A sharp physical-length boundary}
For $a>0$, let $\ell_a(0)=0$ and $\ell_a(t)=[\log(e/t)]^{-a}$ near zero, extended to a continuous
increasing concave function by its tangent at $t=e^{-a}$, and take
the full relation in $\mathbb R^3$
\begin{equation}
 G_a(w_1,w_2,y)=\{(-\ell_a(4y),0,3y),(-\ell_a(4y),0,-5y)\},\quad y\ge0.
\end{equation}
Set $G_a(w)=\varnothing$ for $y<0$.
Its full resolvent is $(p_1,p_2,q)\mapsto
(p_1+\ell_a(|q|),p_2,|q|/4)$; its true residual is
$\chi_a(y)=\sqrt{\ell_a(4y)^2+9y^2}$ at $y\ge0$ and its exact zero-distance gauge
is $\chi_a^{-1}$ on those nonempty fibers. Negative heights have residual
$+\infty$. The reflected modulus
$\sqrt{[t+2\ell_a(t)]^2+9t^2/4}$ is Dini exactly when $a>1$.
For every $\kappa>1/4$, the exact gauge satisfies direct compatibility
on a sufficiently small fixed window. This follows by comparing
$\ell_a(4\kappa t)-\ell_a(t)\ge
a\log(4\kappa)[\log(e/t)]^{-a-1}$ with the $O(t)$ excess in $B(t)$.

On the tangent plane, identify $\mathbb R^2$ with $\mathbb C$ and set
$g(z)=e^{ic/|z|}z$ for $z\ne0$, $g(0)=0$, $c>0$. This is a
homeomorphism, with inverse twist of the opposite sign. Both have a
global modulus $C_c(t+\sqrt t)$, and their largest exponent for a two-point
H\"older bound on any neighborhood of zero is $1/2$. Against zero alone,
both preserve radius and satisfy a Lipschitz bound. To verify the upper bound, let
$\delta=|z-z'|$ and $R=\max(|z|,|z'|)$. If $\delta\ge R/2$,
the output difference is at most $4\delta$. Otherwise it is at most
both $2R$ and $\delta+2c\delta/R$; split at $R=\sqrt\delta$.
For sharpness take radii $r$ and $r/(1+\pi r/c)$ on one ray: their
input difference is of order $r^2$ and their output difference of
order $r$.

Set $v(x_t,y)=(g^{-1}(x_t),y)$ and $F_a=G_a\circ v$. True physical
and kernel zero distances are both $|y|$, and the complete residual is
unchanged, with the finite formula restricted to $y\ge0$. For sufficiently small $y_0>0$, put
\begin{equation}
 y_k=4^{-k}y_0,\quad \rho_k=\sum_{j\ge k}\ell_a(y_j),\quad
 x_k=(g(-\rho_k),y_k).
\end{equation}
For $a>1$, this is a full legitimate GPPA trajectory with kernel
finite length and physical point limit zero. With $b=\log4$ and
$A_k=\log(e/y_0)+kb$,
\begin{equation}
 \rho_k\sim\frac{A_k^{1-a}}{(a-1)b},\quad
 \Delta\theta_k=c(\rho_{k+1}^{-1}-\rho_k^{-1})
 \sim c(a-1)^2b^2 A_k^{a-2}.
\end{equation}
For $1<a<2$, polar chord comparison yields
\begin{equation}
 |g(-\rho_{k+1})-g(-\rho_k)|\sim\frac{c(a-1)}k.
\end{equation}
Physical length is therefore infinite. At $a=2$, choose
$c=\pi/(\log4)^2$; the angle increment tends to $\pi$ and the length
still diverges harmonically. For $a>2$, the bound by
$\rho_k+\rho_{k+1}$ is summable. This calibrated family has the exact
physical-length threshold $a>2$, matching \eqref{eq:physicaldini}
for the half-H\"older inverse. Kernel Dini alone cannot imply
physical finite length under an arbitrary H\"older inverse.

\section{Output errors and the scope of the fusion}
Let $Z=\zer A=v(S)$ be nonempty closed, $R>0$, and
$Q_R=\{z\in Q:d(z,Z)\le R\}$. Suppose the complete kernel iteration
is single-valued and covered on $Q_R$, with
$d(Tz,Z)\le\theta d(z,Z)$, $\norm{Tz-z}\le B(d(z,Z))$, $0<\theta<1$.
For actual output errors $z_{k+1}=Tz_k+e_k\in Q$, $\norm{e_k}\le E_k$,
start from $z_0\in Q_R$ and set $D_0=d(z_0,Z)$. The deterministic envelope is
\begin{equation}
 D_k=\theta^kD_0+\sum_{j<k}\theta^{k-1-j}E_j,
 \quad d(z_k,Z)\le D_k,\quad
 \norm{z_{k+1}-z_k}\le B(D_k)+E_k.
\end{equation}
On a globally available kernel collar, $D_k\le R$ proves each
successive use of coverage. If $E_k\le\delta$ and
$\max(D_0,\delta/(1-\theta))\le R$, then
\begin{equation}
 \limsup_k d(z_k,Z)\le\delta/(1-\theta).
\end{equation}
For the source-style physical output model
$\widehat x_{k+1}\in(v+\lambda F)^{-1}v(x_k)$,
$x_{k+1}=\widehat x_{k+1}+a_k$, a forward modulus
$\norm{v(\widehat x+a)-v(\widehat x)}\le\tau_v(\norm a)$, with
$\tau_v$ nondecreasing, gives
$E_k=\tau_v(\norm{a_k})$ exactly. For a bijective $v$, a global continuous inverse
modulus $\rho$ then gives the physical tube, provided
$\norm{a_k}\le\delta_x$ and
$\max\{D_0,\tau_v(\delta_x)/(1-\theta)\}\le R$,
$\limsup d(x_k,S)\le\rho(\tau_v(\delta_x)/(1-\theta))$.
Persistent errors do not guarantee a point limit: for $Tz=z/2$ and
$e_k=(-1)^k\delta$, the iterates converge to two different parity
limits. The regularized relation $F+\epsilon v$ in source Theorem~4
requires its own certification; it is not automatically covered.

The precise fusion retains the source kernel update, zero-anchor
polarization, and R-continuity bridge, while adding general reflected
moduli, true residual profiles, nonlinear summability, physical lifting,
and sharp failure boundaries. Strictness is established for the fixed
Theorem~2 assumption class and for full-graph same-algorithm
representations. Global novelty, arbitrary lifts or branch reductions,
variable kernels, and the separate global shadow/fiber manuscript are
not silently included in these conclusions.

\begin{thebibliography}{9}
\bibitem{LMT} B. K. Le, B. S. Mordukhovich, and M. A. Th\'era,
\emph{Convergence and Stability Analysis of a Generalized Proximal Point
Algorithm and Its Inexact Version}, arXiv:2608.01584v1, 2026.
\bibitem{BC} M. N. B\`ui and P. L. Combettes,
\emph{Warped proximal iterations for monotone inclusions},
Journal of Mathematical Analysis and Applications 491 (2020), 124315.
This reference supplies terminology and a prior-work boundary; no
unchecked theorem from it is imported into the proofs above.
\end{thebibliography}
\end{document}
